% Generated draft: requires mathematical and source review.
\documentclass[11pt]{article}
\usepackage[margin=25mm]{geometry}
\usepackage{fontspec}
\setmainfont{Latin Modern Roman}
\usepackage{amsmath,amssymb,mathtools,unicode-math}
\setmathfont{Latin Modern Math}
\usepackage{enumitem,needspace,xurl,hyperref}
\setlength{\parindent}{0pt}
\setlength{\parskip}{5pt}
\setlength{\emergencystretch}{3em}
\newcommand{\N}{\mathbb{N}}
\newcommand{\Z}{\mathbb{Z}}
\newcommand{\Q}{\mathbb{Q}}
\newcommand{\R}{\mathbb{R}}
\newcommand{\C}{\mathbb{C}}
\newcommand{\F}{\mathbb{F}}
\newcommand{\A}{\mathbb{A}}
\newcommand{\PP}{\mathbb{P}}
\newcommand{\E}{\mathbb{E}}
\newcommand{\eps}{\varepsilon}
\newcommand{\dd}{\,\mathrm{d}}
\begin{document}
\section*{UnsolvedMath ID: 30003115}
\textbf{Source problem code:} catalogue entry
\subsection{Definitions and mathematical statement}
For a real number \(t\), let \(\{t\}=t-\lfloor t\rfloor\in[0,1)\). For real \(x,y\), positive integers \(a,b\) satisfying \(\gcd(a,b)=1\), and an exponent index set \(I\), define
\[
\mathcal O_I(a,b;x,y)=\bigl\{(\{a^n b^k x\},\{a^n b^k y\}):n,k\in I\bigr\}.
\]
The source writes \(I=\mathbb N\), without specifying whether \(0\in\mathbb N\). Its density conclusion means
\[
\overline{\mathcal O_I(a,b;x,y)}^{\,[0,1]^2}=[0,1]^2,
\]
where closure is taken in the Euclidean topology of the closed square, including its boundary.

\textbf{Source question, with unresolved hypothesis.} The supplied record asks whether, for every pair \(x,y\) satisfying the source condition ``\(1,x,y\) are multiplicatively independent over \(\mathbb Q\)'', and every pair of relatively prime positive integers \(a,b\), the set \(\mathcal O_{\mathbb N}(a,b;x,y)\) is dense in \([0,1]^2\). Real coordinates are required for this interpretation of reduction modulo \(1\), but the source does not explicitly state the domain of \(x,y\). The stated multiplier range includes \(a=1\) or \(b=1\), including \(a=b=1\); no additional lower bound is supplied.

\textbf{Unresolved meaning of independence.} Under the usual definition, elements \(u_1,\ldots,u_r\) of a multiplicative group are multiplicatively independent if
\[
\prod_{j=1}^{r}u_j^{m_j}=1,\quad (m_1,\ldots,m_r)\in\mathbb Z^r
\quad\Longrightarrow\quad m_1=\cdots=m_r=0.
\]
A tuple containing \(1\) cannot satisfy this definition, because the exponent vector \((1,0,0)\) already gives a relation. The qualification ``over \(\mathbb Q\)'' is not defined in the record and does not supply a clear alternative meaning. Consequently the literal hypothesis does not provide a meaningful nonvacuous formulation under the usual definition. Neither replacing multiplicative independence by linear independence over \(\mathbb Q\), nor deleting \(1\) from the tuple, is authorized by the supplied evidence. The intended independence hypothesis, coordinate domain, and multiplier restrictions require mathematical source review before a precise intended conjecture can be stated.
\subsection{Short English statement}
The source asks whether multiplying both coordinates repeatedly by two relatively prime positive integers produces a dense orbit in the unit square. Its independence assumption is unresolved: a tuple containing \(1\) is never multiplicatively independent in the usual sense. The intended hypothesis and the allowed multipliers need review.
\subsection{Sources}
\begin{itemize}
\item Supplied local UnsolvedMath registry record, permanent ID \(30003115\), title ``Dense Multiplicative Orbits on the Two-Torus'', source problem code \texttt{catalogue entry}, in the Oberwolfach Reports Open Problems collection. The record supplies proposed year \(2016\) and names no individual proposer.
\item Supplied source citation: \emph{Diophantische Approximationen} (2017), Oberwolfach Reports, pp.~1099--1187, DOI \texttt{10.4171/owr/2016/21}. Supplied source link: \url{https://doi.org/10.4171/owr/2016/21}. The background describes the reports as documenting mathematical workshops and their open-problem discussions.
\item Supplied literature source, labeled ``OWR report'': \url{https://ems.press/content/serial-article-files/46625?nt=1}.
\item Literature assessment dated exactly \texttt{2026-08-21}: the supplied curation labels the problem open, describes a two-generator Furstenberg density question for multiplicatively independent \(x,y\) as remaining open, and says that known density theorems require stronger hypotheses such as three pairwise coprime multipliers. These are supplied dated curation claims, not independently verified conclusions. The background expressly requests independent verification of time-sensitive status and cited claims; the precise hypotheses of the mentioned density theorems are not supplied.
\item The registry labels the statement \texttt{corrected\_verified}, explaining that mathematical content was preserved while notation, typography, or omitted definitions were repaired. This supplied editorial label does not resolve the independence defect or constitute mathematical verification of this draft. No external source was visited for this preparation.
\end{itemize}
\end{document}
